% Source: OpenAI2026AreaLaw, 06-transport.tex, transport:filtered-vector,
% lines 388--392. Revision adc7f1241b42e322a6451854ab7e4b4c146bf78a.
\section{Filtering under an isometric embedding}
\label{sec:filtering_intertwine}

Let $M$ be positive semidefinite on a finite-dimensional complex Hilbert
space. Define $M^{-1/4}$ by applying $t\mapsto t^{-1/4}$ to its positive
eigenvalues and sending its zero eigenvalues to zero. For a vector $v$, put
\begin{align}
  f_M(v)&=M^{-1/4}v, &
  N_M(v)&=\|f_M(v)\|, &
  V_M(v)&=N_M(v)^{-1}f_M(v),
\end{align}
where the last expression is zero when $N_M(v)=0$.
For positive definite $M$, these are the filtering formulas used in the
transport estimate of~\cite{OpenAI2026AreaLaw}.
The following identities are auxiliary to that construction.

\begin{theorem}[The unnormalized filtered vector under an intertwiner]
  \label{thm:filtered_raw_intertwine}
  \lean{Matrix.Transport.filteredRaw_intertwine}
  \leanok
  Let $M,N$ be positive semidefinite matrices on two finite-dimensional
  complex Hilbert spaces. Suppose that $J$ maps the second space to the
  first and $MJ=JN$. Then every vector $v$ on the second space satisfies
  \begin{align}
    f_M(Jv)&=Jf_N(v).
  \end{align}
\end{theorem}

\begin{proof}
  \leanok
  \uses{thm:real_power_intertwine}
  Apply Theorem~\ref{thm:real_power_intertwine} with exponent $-1/4$
  to obtain $M^{-1/4}J=JN^{-1/4}$, then evaluate this identity at $v$.
\end{proof}

\begin{theorem}[The filtering norm under an isometry]
  \label{thm:filtered_norm_sq_intertwine}
  \lean{Matrix.Transport.filteredNormSq_intertwine}
  \leanok
  Under the preceding hypotheses, suppose also that $J^*J=I$.
  Then every vector $v$ satisfies
  \begin{align}
    N_M(Jv)^2&=N_N(v)^2.
  \end{align}
\end{theorem}

\begin{proof}
  \leanok
  \uses{thm:filtered_raw_intertwine}
  Put $w=f_N(v)$. Theorem~\ref{thm:filtered_raw_intertwine} gives
  $f_M(Jv)=Jw$. The identity $J^*J=I$ then gives
  \begin{align}
    \|Jw\|^2&=\langle Jw,Jw\rangle
       =\langle w,J^*Jw\rangle=\|w\|^2.
  \end{align}
\end{proof}

\begin{theorem}[The normalized filtering formula under an isometry]
  \label{thm:filtered_vector_intertwine}
  \lean{Matrix.Transport.filteredVector_intertwine}
  \leanok
  Under the same hypotheses, every vector $v$ satisfies
  \begin{align}
    V_M(Jv)&=JV_N(v).
  \end{align}
\end{theorem}

\begin{proof}
  \leanok
  \uses{thm:filtered_raw_intertwine, thm:filtered_norm_sq_intertwine}
  Theorem~\ref{thm:filtered_norm_sq_intertwine} gives equality of the
  squared norms and hence of their nonnegative square roots.
  The common scalar $N_N(v)^{-1}$ therefore multiplies both unnormalized vectors.
  Apply Theorem~\ref{thm:filtered_raw_intertwine} and move that scalar
  through $J$. This argument also applies when the common scalar is zero.
\end{proof}

The last identity preserves the normalized formula. A unit-norm conclusion
requires a nonzero filtering norm and is a separate assertion.
